% HD-F1-T15. GPT-6 Astra Pro, /root; 22 November 2026 UTC.
\section[HD-F1-T15: irreducible symplectic-type normal subgroups]{HD-F1-T15: an irreducibly acting symplectic-type normal subgroup excludes the original sixfold source}
\label{sec:f1-symplectic-type-source-exclusion}

The remaining natural-curve gate has no source whose actual coefficient
group has a normal subgroup of symplectic type acting absolutely
irreducibly on the coefficient space. Here the precise group hypothesis
means the following. There is a normal subgroup $R\triangleleft G$
with cyclic centre $Z_R$, with $R/Z_R$ a nonzero elementary abelian
$p$-group, and with the commutator pairing
\[
 R/Z_R\times R/Z_R\longrightarrow (Z_R)[p],\qquad
 (uZ_R,vZ_R)\longmapsto [u,v]
 \tag{F1T15.1}
\]
nondegenerate. The restriction of the original faithful representation
$G\subset GL_K(V)$ to $R$ is absolutely irreducible. This includes
an absolutely irreducibly acting normal extraspecial group and the
specified central products with a cyclic scalar centre. No assertion
about arbitrary repeated constituents of $V|_R$ is included.

For this class the actual inertia matrices in the original source
satisfy
\[
 \operatorname{rank}(g_j-1)\ge\frac m4,
       \qquad m=\dim_K V.
 \tag{F1T15.2}
\]
The original branch budget $\sum_j\operatorname{rank}(g_j-1)=2m+6$
and $b\ge12$ therefore give $m\le6$. This contradicts the remaining
gate $m\ge27$. The proof of (F1T15.2) below concerns the original
matrices and every possible scalar lift.

\paragraph{1. Exact eigenspace dimensions for noncentral elements of $R$.}
Work in the original complexification $V_{\mathbf C}$.
Absolute irreducibility of $R$ makes $Z_R$ scalar; faithfulness
identifies its unique subgroup of order $p$ with the scalar
$p$th roots of unity. If $h\in R\setminus Z_R$, nondegeneracy
of (F1T15.1) supplies $e\in R$ with
\[
                     ehe^{-1}=\zeta h,
 \tag{F1T15.3}
\]
where $\zeta$ is a primitive $p$th root of unity.
Also $h^p\in Z_R$, so the eigenvalues of $h$ are contained in
one scalar coset of the $p$th roots. Equation (F1T15.3) makes
$e$ permute their eigenspaces cyclically. Each of these $p$
eigenspaces consequently has exactly dimension $m/p$.
In particular, for every scalar $t$,
\[
 \operatorname{rank}(h-t)\ge m(1-1/p),\qquad
 \operatorname{rank}(h-1)\ge m(1-1/p).
 \tag{F1T15.4}
\]
This proof permits $h$ to have order $p^2$ or a larger scalar
multiple of order $p$. No order-two replacement is used when $p=2$.

\paragraph{2. The actual commutator map bounds each noncentral eigenspace.}
Let $g\in G$ be nonscalar, and fix any scalar $\lambda$.
Suppose first that some $e\in R$ has
$[g,e]=geg^{-1}e^{-1}\notin Z_R$.
Normality puts this commutator in $R$. The exact matrix identity is
\[
 [g,e]-1=(ge-eg)g^{-1}e^{-1}
       =\bigl((g-\lambda)e-e(g-\lambda)\bigr)g^{-1}e^{-1}.
 \tag{F1T15.5}
\]
Taking ranks of these actual maps and using (F1T15.4) gives
\[
 m(1-1/p)\le\operatorname{rank}([g,e]-1)
                    \le2\operatorname{rank}(g-\lambda).
 \tag{F1T15.6}
\]
Hence every eigenspace of $g$ has codimension at least
$m(p-1)/(2p)\ge m/4$ in this case.

In the other case $[g,e]\in Z_R$ for every $e\in R$.
Since $Z_R$ is scalar, conjugation by $g$ fixes it pointwise.
The rule
\[
 \chi_g(eZ_R)=geg^{-1}e^{-1}
 \tag{F1T15.7}
\]
is therefore a homomorphism from $R/Z_R$ to $Z_R$.
Its values have order dividing $p$: conjugation fixes $e^p\in Z_R$,
so $\chi_g(eZ_R)^p=1$. The commutator pairing gives an injective
homomorphism
\[
 R/Z_R\longrightarrow\operatorname{Hom}(R/Z_R,(Z_R)[p]),
       \qquad uZ_R\longmapsto\bigl(eZ_R\mapsto[u,e]\bigr).
 \tag{F1T15.8}
\]
Both finite groups have the same order, and thus this displayed
map is an isomorphism. Choose the resulting actual $u\in R$
with $[u,e]=\chi_g(eZ_R)$ for every $e$.
Then $u^{-1}g$ centralizes $R$, and absolute irreducibility
makes it a scalar. Thus $g=zu$ for a scalar $z$.
Because $g$ is nonscalar, $u\notin Z_R$; (F1T15.4) proves
the stronger eigenspace codimension bound $m(1-1/p)$.

These two cases prove the bound $m/4$ for every eigenspace
of every nonscalar $g\in G$. A nonidentity scalar inertia
matrix has $\operatorname{rank}(g-1)=m$ directly.
This proves (F1T15.2) for every actual nonidentity inertia.

\paragraph{3. The original branch formula and the retained polarized factor.}
For the same original cover, the middle-extension stalk at the
$j$th branch value is $V^{g_j}$. Its parabolic cohomology has
dimension six, and the period image of dimension nine requires
at least twelve true branch values, as in (F1T13.1). Therefore
\[
       2m+6=\sum_{j=1}^b\operatorname{rank}(g_j-1)
                    \ge\frac{bm}{4}\ge3m,
                  \qquad m\le6.
 \tag{F1T15.9}
\]
In particular this whole source class is absent in the remaining
moving-rank-at-least-five branch.

The displayed morphisms are conjugation within the actual coefficient
group, the commutator map (F1T15.5), and the perfect finite pairing
(F1T15.8). The subgroup $R$ is used only to bound ranks inside the
original $V$. The cover, middle-extension sheaf, rational primitive
idempotent in $\mathbf Q[G]$, actual primitive Jacobian image, field
action and restricted theta polarization are retained by their identity
maps. Their parabolic cup product and Poincare biextension maps are
also identities, with multiplier one. The same identity realizes this
exclusion in the variational, Grothendieck-group, Adams and weight-one
frameworks of F01--F04. No new determinant component or effective cycle
is supplied by the exclusion.

The elementary finite pairing and matrix-rank proof above are complete;
no classification of primitive linear groups is a premise. The term
``symplectic type'' is used only with the exact hypotheses preceding
(F1T15.1). A normal subgroup with several irreducible constituents, or
with a nontrivial multiplicity space, is left at its existing gate.
